\documentclass[10pt,a4paper]{amsart}
\usepackage[margin=19mm]{geometry}
\usepackage{amsmath,amssymb,mathtools}
\usepackage[scr=rsfs,cal=euler]{mathalfa}
\usepackage{xfrac}
\usepackage[unicode,hidelinks,bookmarks=false]{hyperref}
\newcommand{\ie}{i.e.\ }
\newcommand{\cf}{cf.\ }
\newcommand{\resp}{respectively\ }
\newcommand{\Subsection}{\subsection}
\providecommand{\nobreakdash}{\nobreak\hbox{-}}
\setlength{\emergencystretch}{2em}
\multlinegap0pt
\usepackage{amssymb}
\newtheorem{lemma}{Lemma}
\newtheorem{prop}{Proposition}
\newcommand{\floor}[1]{\left \lfloor #1 \right \rfloor}
\title{A determinant of the Artin-Hasse exponential coefficients}
\author{Matthew Schmidt}
\date{Source: arXiv:2106.06084v1 (CC BY 4.0)}
\begin{document}
\maketitle

\noindent\emph{Source and licence.} A determinant of the Artin-Hasse exponential coefficients. Version arXiv:2106.06084v1 (CC BY 4.0), distributed by arXiv under the Creative Commons Attribution 4.0 International licence, http://creativecommons.org/licenses/by/4.0/, as recorded in the arXiv licence metadata for this exact version and shown on the abstract page https://arxiv.org/abs/2106.06084v1. Full licence notice: \texttt{licenses/arxiv-2106.06084-CC-BY-4.0.txt}.\par\smallskip

\noindent\textit{Editorial scope.} Counted unit: the article's prop prop:CD\_det (source0.tex lines 377--382). The article's complete original proof of it is delivered with it (source0.tex lines 383--409, one \texttt{proof} environment of 27 lines). No mathematics is written, repaired or re-derived: the statement and the proof are the article's own lines, in the article's own notation, and no retained line is substituted or re-flowed.  Retained uncounted as the source-local context the proof needs: lines 337--343; lines 359--364.  Nothing was omitted from any retained span, so the package carries no omission record.  No retained line contains a citation, so the wrapper carries no bibliography and no bibliography entry.  No retained line contains a figure environment, an image-inclusion command or drawing code, so no image is delivered and the package declares no asset dependency.  Editorial formatting only, in addition: an AMS \texttt{amsart} preamble that restates the article's own declarations and macros used by the retained lines, namely the environments \texttt{lemma}, \texttt{prop} on separate counters, together with the standard packages those lines need.  The retained environments print in this document as Lemma 1, Proposition 1 in order of appearance, whereas the article prints the same items with the numbering of its own full document.  The complete native source of the article as distributed in the arXiv e-print for this exact version is bundled unchanged as \texttt{sources/112528/source0.tex}.
\begin{lemma}\label{lemma:h_n_expans}
For $n\geq 1$,
\begin{align*}
	h_n = \sum_{k=0}^{\floor{n/p}} \binom{n}{kp} C_k,
\end{align*}	
where $C_k = \sum_{t\in I^k}C_t^{kp}$.
\end{lemma}

\begin{lemma}\label{lemma:deter}
For any formal variables $E_{j,k}$ and $X_{i,k}$, 
\[
	|  \sum_{k=1}^{\ell} E_{j, i-k+1} X_{i,k} |_{1\leq i,j\leq \ell} = |E_{j,k}|_{1\leq j,k\leq \ell} \cdot \prod_{i=1}^\ell X_{i,1}.
\]
\end{lemma}

\begin{prop}\label{prop:CD_det}
Let $i,j\geq 1$ and $\alpha_j,\beta_i\in\mathbb{Q}_{>0}$. For any $\ell\geq 1$,
\begin{align*}
	|\alpha_j\beta_iu_{pi-j}|_{1\leq i,j\leq \ell} =  |\frac{\alpha_j\beta_i}{(pk-j)!}|_{1\leq j,k\leq \ell}.
\end{align*}
\end{prop}

\begin{proof}
Let $i,j\geq 1$ with $pi-j\geq 0$. By Lemma~\ref{lemma:h_n_expans}:
\begin{align*}
	h_{pi-j} = \sum_{k=1}^{i-\floor{j/p}}\binom{ip-j}{kp} C_k.
\end{align*}
Hence:
\begin{align}\label{line:CjDih}
	\frac{\alpha_j\beta_i}{(pi-j)!}h_{pi-j} &= \sum_{k=0}^{i-\floor{j/p}}\left (\frac{\alpha_j}{(p(i-k)-j)!}\right )\left (\beta_i\frac{C_k}{(pk)!}\right).
\end{align}

Thus if $X_{i,k} = \beta_i\frac{C_k}{(pk)!}$ and
\begin{align*}
E_{j, k} = \begin{cases}
		\frac{\alpha_j}{(pk-j)!} & k\geq  \floor{j/p}\\
		0 & k< \floor{j/p}
		\end{cases},
\end{align*} 
then Equation~(\ref{line:CjDih}) becomes:
\begin{align}
\frac{\alpha_j\beta_i}{(pi-j)!}h_{pi-j} &= \sum_{k=0}^{\ell}E_{j,i-k}X_{i,k}.
\end{align}
So by Lemma~\ref{lemma:deter},
\begin{align*}
	|\alpha_j\beta_iu_{pi-j}|_{1\leq i,j\leq \ell} &=|\frac{\alpha_j\beta_i}{(pi-j)!}h_{pi-j}|_{1\leq i,j\leq \ell} =  |E_{j,k}|_{1\leq j,k\leq \ell} \cdot \prod_{i=1}^\ell X_{i,0} \\
	&=  |\frac{\alpha_j\beta_k}{(pk-j)!}|_{1\leq j,k\leq \ell}.
\end{align*}
\end{proof}

\end{document}
